% GENERATED FILE. Edit canonical lessons, then run npm run generate:books.
% canonical-source-hash: fnv1a64:a77a66e10f70aa68
% canonical-lessons: KA-C201-mundudu, KA-C201-vivarisu, KA-C201-pasagu, KA-C201-nenapidu, KA-C201-mudrisu

\chapter{To postpone, To explain, To pass, To remember, To print}
\label{ch:ka-to-postpone-to-explain-to-pass-to-remember-to-print}
\hlchaptermodality{201}

\begin{chapteropening}
\textbf{By the end of this chapter:} \emph{I can say to postpone, to explain, to pass, to remember and to print.}

Complete the last lesson of chapter 201: I can say to postpone, to explain, to pass, to remember and to print.
\end{chapteropening}

\section[\texorpdfstring{\kn{ಮುಂದೂಡು} (mundūḍu)}{ಮುಂದೂಡು (mundūḍu)}]{\kn{ಮುಂದೂಡು} (mundūḍu) — to postpone}

\label{lesson:KA-C201-mundudu}



\begin{quote}
Before the new one: say the Kannada for to finish, then the Kannada for to start.
\end{quote}

\subsection*{You'll want to know: \kn{ಮುಂದೂಡು}}
\textbf{\kn{ಮುಂದೂಡು}} — \emph{mundūḍu} — \textquotedblleft{}to postpone\textquotedblright{}.

\begin{grammarlens}[title={What you've built}]
One more word for this chapter.
\end{grammarlens}

\subsection*{Your turn}
\begin{itemize}
\raggedright
  \item \emph{Say it:} \emph{mundūḍu}
  \item \emph{Say it:} \emph{mundūḍu}, once more
  \item \emph{Say it:} say \emph{mundūḍu}
  \item \emph{Recall:} say the Kannada for to finish, then the Kannada for to start, then say \emph{mundūḍu} again
\end{itemize}

\begin{tcolorbox}[breakable,colback=teal!4,colframe=teal!35!black,title={Before you move on}]
What does \kn{ಮುಂದೂಡು} mean? (\textquotedblleft{}To postpone\textquotedblright{}.) Say it once more.
\end{tcolorbox}
\section[\texorpdfstring{\kn{ವಿವರಿಸು} (vivarisu)}{ವಿವರಿಸು (vivarisu)}]{\kn{ವಿವರಿಸು} (vivarisu) — to explain}

\label{lesson:KA-C201-vivarisu}



\begin{quote}
Before the new one: say the Kannada for to start, then the Kannada for to postpone.
\end{quote}

\subsection*{You'll want to know: \kn{ವಿವರಿಸು}}
\textbf{\kn{ವಿವರಿಸು}} — \emph{vivarisu} — \textquotedblleft{}to explain\textquotedblright{}.

\begin{grammarlens}[title={What you've built}]
One more word for this chapter.
\end{grammarlens}

\subsection*{Your turn}
\begin{itemize}
\raggedright
  \item \emph{Say it:} \emph{vivarisu}
  \item \emph{Say it:} \emph{vivarisu}, once more
  \item \emph{Say it:} say \emph{vivarisu}
  \item \emph{Recall:} say the Kannada for to start, then the Kannada for to postpone, then say \emph{vivarisu} again
\end{itemize}

\begin{tcolorbox}[breakable,colback=teal!4,colframe=teal!35!black,title={Before you move on}]
What does \kn{ವಿವರಿಸು} mean? (\textquotedblleft{}To explain\textquotedblright{}.) Say it once more.
\end{tcolorbox}
\section[\texorpdfstring{\kn{ಪಾಸಾಗು} (pāsāgu)}{ಪಾಸಾಗು (pāsāgu)}]{\kn{ಪಾಸಾಗು} (pāsāgu) — to pass (an exam)}

\label{lesson:KA-C201-pasagu}



\begin{quote}
Before the new one: say the Kannada for to postpone, then the Kannada for to explain.
\end{quote}

\subsection*{You'll want to know: \kn{ಪಾಸಾಗು}}
\textbf{\kn{ಪಾಸಾಗು}} — \emph{pāsāgu} — \textquotedblleft{}to pass (an exam)\textquotedblright{}.

\begin{grammarlens}[title={What you've built}]
One more word for this chapter.
\end{grammarlens}

\subsection*{Your turn}
\begin{itemize}
\raggedright
  \item \emph{Say it:} \emph{pāsāgu}
  \item \emph{Say it:} \emph{pāsāgu}, once more
  \item \emph{Say it:} say \emph{pāsāgu}
  \item \emph{Recall:} say the Kannada for to postpone, then the Kannada for to explain, then say \emph{pāsāgu} again
\end{itemize}

\begin{tcolorbox}[breakable,colback=teal!4,colframe=teal!35!black,title={Before you move on}]
What does \kn{ಪಾಸಾಗು} mean? (\textquotedblleft{}To pass (an exam)\textquotedblright{}.) Say it once more.
\end{tcolorbox}
\section[\texorpdfstring{\kn{ನೆನಪಿಡು} (nenapiḍu)}{ನೆನಪಿಡು (nenapiḍu)}]{\kn{ನೆನಪಿಡು} (nenapiḍu) — to remember, to keep in mind}

\label{lesson:KA-C201-nenapidu}



\begin{quote}
Before the new one: say the Kannada for to explain, then the Kannada for to pass.
\end{quote}

\subsection*{You'll want to know: \kn{ನೆನಪಿಡು}}
\textbf{\kn{ನೆನಪಿಡು}} — \emph{nenapiḍu} — \textquotedblleft{}to remember, to keep in mind\textquotedblright{}.

\begin{grammarlens}[title={What you've built}]
One more word for this chapter.
\end{grammarlens}

\subsection*{Your turn}
\begin{itemize}
\raggedright
  \item \emph{Say it:} \emph{nenapiḍu}
  \item \emph{Say it:} \emph{nenapiḍu}, once more
  \item \emph{Say it:} say \emph{nenapiḍu}
  \item \emph{Recall:} say the Kannada for to explain, then the Kannada for to pass, then say \emph{nenapiḍu} again
\end{itemize}

\begin{tcolorbox}[breakable,colback=teal!4,colframe=teal!35!black,title={Before you move on}]
What does \kn{ನೆನಪಿಡು} mean? (\textquotedblleft{}To remember, to keep in mind\textquotedblright{}.) Say it once more.
\end{tcolorbox}
\section[\texorpdfstring{\kn{ಮುದ್ರಿಸು} (mudrisu)}{ಮುದ್ರಿಸು (mudrisu)}]{\kn{ಮುದ್ರಿಸು} (mudrisu) — to print}

\label{lesson:KA-C201-mudrisu}



\begin{quote}
Before the new one: say the Kannada for to pass, then the Kannada for to remember.
\end{quote}

\subsection*{You'll want to know: \kn{ಮುದ್ರಿಸು}}
\textbf{\kn{ಮುದ್ರಿಸು}} — \emph{mudrisu} — \textquotedblleft{}to print\textquotedblright{}.

\begin{grammarlens}[title={What you've built}]
One more word for this chapter.
\end{grammarlens}

\subsection*{Your turn}
\begin{itemize}
\raggedright
  \item \emph{Say it:} \emph{mudrisu}
  \item \emph{Say it:} \emph{mudrisu}, once more
  \item \emph{Say it:} say \emph{mudrisu}
  \item \emph{Recall:} say the Kannada for to pass, then the Kannada for to remember, then say \emph{mudrisu} again
\end{itemize}

\begin{tcolorbox}[breakable,colback=teal!4,colframe=teal!35!black,title={Before you move on}]
What does \kn{ಮುದ್ರಿಸು} mean? (\textquotedblleft{}To print\textquotedblright{}.) Say it once more.
\end{tcolorbox}
